% =========================
% Explanation: Algorithms Compared (informal, readable)
% =========================
\subsection*{Explanation of Algorithms Compared (informal)}
This section clarifies \emph{which optimizers we compare}, \emph{why each baseline is included}, and \emph{how we guarantee a fair comparison}.
Clarity is essential in optimization benchmarking: reviewers must be able to tell exactly what information and resources each method uses.

\begin{itemize}
    \item \textbf{Proposed method as an input--output mapping.}
    We treat each optimizer as a procedure that maps past observations to new configurations.
    At step \(t\), the optimizer has access to a history
    \[
        \mathcal{H}_t = \{(\lambda_i, y_i)\}_{i=1}^{t}
        \quad \text{(or vector outcomes \(\mathbf{y}_i\) in multiobjective settings),}
    \]
    and outputs one or more new candidates \(\lambda_{t+1}\) (or a batch) to evaluate next.
    A reproducible protocol must specify:
    \begin{enumerate}
        \item what the optimizer conditions on (full history vs.\ an elite subset),
        \item how candidates are generated (sampling / mutation / crossover / acquisition),
        \item how solutions are selected or archived (incumbent best, nondominated archive, etc.),
        \item how invalid configurations are handled (repair, rejection, or penalty).
    \end{enumerate}

    \item \textbf{What ``P3Net'' means in this paper.}
    In our terminology, \textbf{P3Net} means:
    \begin{itemize}
        \item \textbf{P3 is the backbone:} the optimizer fits a probabilistic generator \(q_{\theta}(\lambda)\) to an elite subset of evaluated configurations
        and proposes new candidates by sampling from \(q_{\theta}\).
        \item \textbf{``Net'' is the adaptation layer:} additional mechanisms required to apply P3 fairly and effectively to the benchmarked HPO/NAS setting:
        \begin{itemize}
            \item a canonical encoding of \(\lambda\) (vector or genotype) used by all methods;
            \item an \emph{elite selection} rule specifying which evaluated configurations are used to fit \(q_{\theta}\);
            \item a \emph{diversity / anti-collapse} mechanism to avoid degenerating to a single mode;
            \item \emph{constraint-aware sampling} (masking/repair/rejection) so samples are valid under conditionals or NAS constraints;
            \item a \emph{batched proposal} policy compatible with parallel evaluation (propose up to \(W\) configurations at a time);
            \item an \emph{update schedule} specifying when \(q_{\theta}\) is refit and how multi-fidelity information (if any) is incorporated.
        \end{itemize}
    \end{itemize}
    All of these components must be fixed \emph{a priori} for the benchmark suite to avoid post-hoc tuning.

    \item \textbf{Baselines and why they are appropriate.}
    We include baselines spanning the major families of HPO/NAS methods:
    \begin{itemize}
        \item \textbf{Random Search (RS):} a required reference point for sample efficiency; often strong in high dimensions.
        \item \textbf{Model-based HPO (BO/TPE/SMAC):} standard, strong baselines for expensive black-box objectives.
        TPE and SMAC are especially appropriate for mixed and conditional search spaces common in HPO/NAS.
        \item \textbf{CMA-ES / evolutionary strategies (when applicable):} competitive derivative-free optimization for continuous spaces,
        serving as a strong non-model-based baseline.
        \item \textbf{Multi-fidelity schedulers (Hyperband/ASHA, BOHB) when allowed:} if evaluations admit fidelities
        (e.g., epochs/tokens/resolution), these baselines allocate resources adaptively and are highly competitive.
        \item \textbf{Evolutionary Pareto baselines (e.g., NSGA-Net / NSGA-II) when multiobjective:} appropriate comparators for Pareto-front approximation
        under the same encoding and evaluation model.
        \item \textbf{Component attribution baselines:} \emph{vanilla P3} (without ``Net'' adaptations) isolates the contribution of P3Net mechanisms
        (elite rule, diversity, constraint handling, batching, update schedule).
    \end{itemize}

    \item \textbf{Baselines we cannot run and feasibility constraints.}
    Not all published methods are directly comparable under a fixed evaluation model.
    For each omitted baseline, we give a concrete reason, e.g.:
    \begin{itemize}
        \item incompatible evaluation model (e.g., weight-sharing NAS changes what ``one evaluation'' means),
        \item missing or unsupported encoding requirements (e.g., specialized graph kernels),
        \item prohibitive compute or memory costs under our parallelism/budget constraints,
        \item or implementation / licensing unavailability.
    \end{itemize}
    We narrow claims accordingly: we only claim improvements over methods that can be evaluated fairly under the same protocol.

    \item \textbf{Fair comparison rules (hard constraints).}
    To guarantee fairness, we enforce:
    \begin{enumerate}
        \item \textbf{Identical search space \(\Lambda\) and encoding} for all methods.
        \item \textbf{Identical evaluation pipeline} (data splits, preprocessing, training recipe, metric computation).
        \item \textbf{Identical budgets and stopping rules:} each method receives the same evaluation budget \(B\);
        invalid/failed trials are handled consistently and count toward the budget.
        \item \textbf{Identical parallelism policy:} same worker count \(W\) and identical synchronous/asynchronous dispatcher semantics.
        \item \textbf{Identical randomization control:} same seed policy and, when used, common random numbers for evaluation stochasticity.
    \end{enumerate}
    These constraints ensure differences reflect the optimization strategy, not extra information or resources.

    \item \textbf{What comparisons are feasible.}
    Comparisons are feasible when methods (i) operate on the same \(\Lambda\), (ii) use the same evaluation unit and objective definition,
    and (iii) follow the same budget accounting.
    If any of these conditions are violated (e.g., a method changes the evaluation model), we treat it as \emph{not directly comparable}
    and exclude it or evaluate it separately under a clearly labeled protocol.
\end{itemize}
